% Options for packages loaded elsewhere
\PassOptionsToPackage{unicode}{hyperref}
\PassOptionsToPackage{hyphens}{url}
\PassOptionsToPackage{dvipsnames,svgnames,x11names}{xcolor}
%
\documentclass[
  letterpaper,
  DIV=11,
  numbers=noendperiod]{scrartcl}

\usepackage{amsmath,amssymb}
\usepackage{iftex}
\ifPDFTeX
  \usepackage[T1]{fontenc}
  \usepackage[utf8]{inputenc}
  \usepackage{textcomp} % provide euro and other symbols
\else % if luatex or xetex
  \usepackage{unicode-math}
  \defaultfontfeatures{Scale=MatchLowercase}
  \defaultfontfeatures[\rmfamily]{Ligatures=TeX,Scale=1}
\fi
\usepackage{lmodern}
\ifPDFTeX\else  
    % xetex/luatex font selection
\fi
% Use upquote if available, for straight quotes in verbatim environments
\IfFileExists{upquote.sty}{\usepackage{upquote}}{}
\IfFileExists{microtype.sty}{% use microtype if available
  \usepackage[]{microtype}
  \UseMicrotypeSet[protrusion]{basicmath} % disable protrusion for tt fonts
}{}
\makeatletter
\@ifundefined{KOMAClassName}{% if non-KOMA class
  \IfFileExists{parskip.sty}{%
    \usepackage{parskip}
  }{% else
    \setlength{\parindent}{0pt}
    \setlength{\parskip}{6pt plus 2pt minus 1pt}}
}{% if KOMA class
  \KOMAoptions{parskip=half}}
\makeatother
\usepackage{xcolor}
\setlength{\emergencystretch}{3em} % prevent overfull lines
\setcounter{secnumdepth}{5}
% Make \paragraph and \subparagraph free-standing
\ifx\paragraph\undefined\else
  \let\oldparagraph\paragraph
  \renewcommand{\paragraph}[1]{\oldparagraph{#1}\mbox{}}
\fi
\ifx\subparagraph\undefined\else
  \let\oldsubparagraph\subparagraph
  \renewcommand{\subparagraph}[1]{\oldsubparagraph{#1}\mbox{}}
\fi


\providecommand{\tightlist}{%
  \setlength{\itemsep}{0pt}\setlength{\parskip}{0pt}}\usepackage{longtable,booktabs,array}
\usepackage{calc} % for calculating minipage widths
% Correct order of tables after \paragraph or \subparagraph
\usepackage{etoolbox}
\makeatletter
\patchcmd\longtable{\par}{\if@noskipsec\mbox{}\fi\par}{}{}
\makeatother
% Allow footnotes in longtable head/foot
\IfFileExists{footnotehyper.sty}{\usepackage{footnotehyper}}{\usepackage{footnote}}
\makesavenoteenv{longtable}
\usepackage{graphicx}
\makeatletter
\def\maxwidth{\ifdim\Gin@nat@width>\linewidth\linewidth\else\Gin@nat@width\fi}
\def\maxheight{\ifdim\Gin@nat@height>\textheight\textheight\else\Gin@nat@height\fi}
\makeatother
% Scale images if necessary, so that they will not overflow the page
% margins by default, and it is still possible to overwrite the defaults
% using explicit options in \includegraphics[width, height, ...]{}
\setkeys{Gin}{width=\maxwidth,height=\maxheight,keepaspectratio}
% Set default figure placement to htbp
\makeatletter
\def\fps@figure{htbp}
\makeatother

\usepackage{booktabs}
\usepackage{longtable}
\usepackage{array}
\usepackage{multirow}
\usepackage{wrapfig}
\usepackage{float}
\usepackage{colortbl}
\usepackage{pdflscape}
\usepackage{tabu}
\usepackage{threeparttable}
\usepackage{threeparttablex}
\usepackage[normalem]{ulem}
\usepackage{makecell}
\usepackage{xcolor}
\KOMAoption{captions}{tableheading}
\makeatletter
\makeatother
\makeatletter
\makeatother
\makeatletter
\@ifpackageloaded{caption}{}{\usepackage{caption}}
\AtBeginDocument{%
\ifdefined\contentsname
  \renewcommand*\contentsname{Table of contents}
\else
  \newcommand\contentsname{Table of contents}
\fi
\ifdefined\listfigurename
  \renewcommand*\listfigurename{List of Figures}
\else
  \newcommand\listfigurename{List of Figures}
\fi
\ifdefined\listtablename
  \renewcommand*\listtablename{List of Tables}
\else
  \newcommand\listtablename{List of Tables}
\fi
\ifdefined\figurename
  \renewcommand*\figurename{Figure}
\else
  \newcommand\figurename{Figure}
\fi
\ifdefined\tablename
  \renewcommand*\tablename{Table}
\else
  \newcommand\tablename{Table}
\fi
}
\@ifpackageloaded{float}{}{\usepackage{float}}
\floatstyle{ruled}
\@ifundefined{c@chapter}{\newfloat{codelisting}{h}{lop}}{\newfloat{codelisting}{h}{lop}[chapter]}
\floatname{codelisting}{Listing}
\newcommand*\listoflistings{\listof{codelisting}{List of Listings}}
\makeatother
\makeatletter
\@ifpackageloaded{caption}{}{\usepackage{caption}}
\@ifpackageloaded{subcaption}{}{\usepackage{subcaption}}
\makeatother
\makeatletter
\@ifpackageloaded{tcolorbox}{}{\usepackage[skins,breakable]{tcolorbox}}
\makeatother
\makeatletter
\@ifundefined{shadecolor}{\definecolor{shadecolor}{rgb}{.97, .97, .97}}
\makeatother
\makeatletter
\makeatother
\makeatletter
\makeatother
\ifLuaTeX
  \usepackage{selnolig}  % disable illegal ligatures
\fi
\IfFileExists{bookmark.sty}{\usepackage{bookmark}}{\usepackage{hyperref}}
\IfFileExists{xurl.sty}{\usepackage{xurl}}{} % add URL line breaks if available
\urlstyle{same} % disable monospaced font for URLs
\hypersetup{
  pdftitle={Feature Engineering and Conformal Prediction in Fish Classification.},
  pdfauthor={Mariana Garcia Mejia},
  colorlinks=true,
  linkcolor={blue},
  filecolor={Maroon},
  citecolor={Blue},
  urlcolor={Blue},
  pdfcreator={LaTeX via pandoc}}

\title{Feature Engineering and Conformal Prediction in Fish
Classification.\thanks{Code and data are available at:
\url{https://github.com/mejiama16/Fish-Project}. The data for this paper
was retrieved from the
\href{https://github.com/WidebandPingFest/FishTetherExperiment}{Fish
Tether Experiment GitHub Repository}.}}
\author{Mariana Garcia Mejia}
\date{October 8, 2026}

\begin{document}
\maketitle
\begin{abstract}
Description of methods and main findings
\end{abstract}
\ifdefined\Shaded\renewenvironment{Shaded}{\begin{tcolorbox}[boxrule=0pt, breakable, frame hidden, borderline west={3pt}{0pt}{shadecolor}, enhanced, interior hidden, sharp corners]}{\end{tcolorbox}}\fi

\hypertarget{sec-introduction}{%
\section{Introduction}\label{sec-introduction}}

Hydroacoustic signals provide a non-lethal approach for fish
classification. However, the features used for classification are
high-dimensional and highly correlated. We are also often only given
labels for a subset of the data, and must leverage both labelled and
non-labelled data to improve classification accuracy. Lastly, we often
care to understand which signals are most important for classification,
but the correlation among features distorts this signal. The student
will use various data science and machine learning approaches to address
the data structure challenges and classification tasks.

\hypertarget{sec-data}{%
\section{Data}\label{sec-data}}

The data set contains \textbf{14,575 rows} and 484 columns. Each row
represents one fish and it's measurement for a single ``ping''. A
``ping'' refers to a time the fish hydroacoustic signal was measured and
stored in the dataset. Since multiple signals were recorded for a single
fish, there are several rows corresponding the same fish (e.g.~there are
multiple rows for LT001; Lake Trout fish number one, because multiple
measures were taken for this single fish). From now on, we call rows
``pings'', to familiarize ourselves with how the data was being
collected.

The dataset contains information for three fish species: Lake Trout
(LT), Lake Whitefish (LW) and Smallmouth Bass (SMB).
Table~\ref{tbl-data-overview} summarizes relevant information about the
dataset for each of the species.

\hypertarget{tbl-data-overview}{}
\begin{longtable}[]{@{}
  >{\raggedright\arraybackslash}p{(\columnwidth - 10\tabcolsep) * \real{0.1067}}
  >{\centering\arraybackslash}p{(\columnwidth - 10\tabcolsep) * \real{0.2133}}
  >{\centering\arraybackslash}p{(\columnwidth - 10\tabcolsep) * \real{0.1733}}
  >{\centering\arraybackslash}p{(\columnwidth - 10\tabcolsep) * \real{0.1867}}
  >{\centering\arraybackslash}p{(\columnwidth - 10\tabcolsep) * \real{0.1600}}
  >{\centering\arraybackslash}p{(\columnwidth - 10\tabcolsep) * \real{0.1600}}@{}}
\caption{\label{tbl-data-overview}Overview of the
dataset}\tabularnewline
\toprule\noalign{}
\begin{minipage}[b]{\linewidth}\raggedright
Species
\end{minipage} & \begin{minipage}[b]{\linewidth}\centering
Number of fish
\end{minipage} & \begin{minipage}[b]{\linewidth}\centering
Total pings
\end{minipage} & \begin{minipage}[b]{\linewidth}\centering
Median pings
\end{minipage} & \begin{minipage}[b]{\linewidth}\centering
Min. pings
\end{minipage} & \begin{minipage}[b]{\linewidth}\centering
Max. pings
\end{minipage} \\
\midrule\noalign{}
\endfirsthead
\toprule\noalign{}
\begin{minipage}[b]{\linewidth}\raggedright
Species
\end{minipage} & \begin{minipage}[b]{\linewidth}\centering
Number of fish
\end{minipage} & \begin{minipage}[b]{\linewidth}\centering
Total pings
\end{minipage} & \begin{minipage}[b]{\linewidth}\centering
Median pings
\end{minipage} & \begin{minipage}[b]{\linewidth}\centering
Min. pings
\end{minipage} & \begin{minipage}[b]{\linewidth}\centering
Max. pings
\end{minipage} \\
\midrule\noalign{}
\endhead
\bottomrule\noalign{}
\endlastfoot
LT & 21 & 7,248 & 213 & 16 & 1,379 \\
LW & 11 & 4,285 & 310 & 33 & 1,619 \\
SMB & 12 & 3,042 & 210 & 14 & 1,071 \\
\end{longtable}

The data was obtained from 44 distinct fishes: 21 Lake Trout, 11 Lake
Whitefish and and 12 Smallmouth Bass. Lake Trout fishes has more pings
(instances where data was collected from them), followed by Lake
Whitefish and Smallmouth Bass. Although the LW species had less
individual fish, more pings were obtained for this species than for SMB.
The median pings column shows the median number of pings across all
individual fishes for a certain species, in other words, what was the
median number of instances where measurements were taken from a single
fish. The values are fairly similar across LT and SMB, while LW had more
pings in a single fish compared to the other. This is likely because the
minimum and maximum number of pings obtained for LW is much higher than
for LT and SMB.

It is important to understand the structure of the data and how it is
collected, as it largely shapes the statistical method we select and how
we make decisions throughout the process. In particular, the rows in
this data do not represent single instances or observations, rather,
each represents an instance of a subject in time; a measurement taken
from a fish. As a result, it is possible for measurements obtained in
these instances, although different, may be highly correlated as they
come from the same fish. For instance, the rows 1-65 in the raw data set
correspond to LT001, Lake Trout fish number one. Because all
hydroacoustic signals are taken from this single fish, just at different
points in time, it is possible for these to vary very little from
measurement to measurement, resulting in highly correlated observations.
Additionally, this feature of the data likely violates the assumption of
independence of observations, that must hold from the simplest models
(OLS regression) to the most complex models in machine learning. We are
aware of this violation and consider its implications accordingly.

Now, one might think that the solution is to constraint our data set and
analysis to just the 44 individual fish, by randomly selecting one of
the measurements taken from each of them. The first issue with this
approach is the resulting sample size. We have over 14,000 measurements
obtained from this experiment, it seems like exploring the data to
overcome the high correlation issue would be a much better solution, as
compared to blindly throwing observations away. Another issue with this
is: how do we pick which individual pings to keep and which to ignore?
Do we hand-pick them? Do we select them randomly? Do we select two per
day and ignore the rest? Lastly, when building machine learning models
with this data, how can we ensure there is no data leakage? If we keep
some observations from fish LT001 in the training set and some in the
testing set, it is possible that due to the high correlation between
these, the model has already memorized these patterns, leading to an
overestimation of the model's actual capacity.

This paper will touch on some of the decisions we must make when
building models using this kind of data, and compare across them. We
will do so in the context of dimensionality reduction and feature
engineering. In particular, we want to get information regarding which
features are the most informative in the classification task. That is,
given all the signals we have available for a new, unknown fish, which
of these signals are the most likely to help us decide whether this fish
is a Lake Trout, a Lake Whitefish or a Smallmouth Bass.

\hypertarget{time-series-component}{%
\subsection{Time Series Component}\label{time-series-component}}

Before jumping into the methods it is important to better understand
other characteristics of the data. First, the data contains a time
series component. The pings for each fish are taken at different times
throughout a period of time, no two pings happen simultaneously.
Table~\ref{tbl-data-time} shows that the data was collected in the span
of a week, with SMB having data collected in a much shorter period of
time. Note that the data set included other dates, for example
\texttt{processedDate}, however we keep \texttt{sampleDate} to reflect
the time the measurement was taken.

\hypertarget{tbl-data-time}{}
\begin{longtable}[]{@{}lcc@{}}
\caption{\label{tbl-data-time}Span of the data collection process
(pings)}\tabularnewline
\toprule\noalign{}
Species & Earliest day a ping was made & Latest day a ping was made \\
\midrule\noalign{}
\endfirsthead
\toprule\noalign{}
Species & Earliest day a ping was made & Latest day a ping was made \\
\midrule\noalign{}
\endhead
\bottomrule\noalign{}
\endlastfoot
LT & 2022-07-21 & 2022-07-28 \\
LW & 2022-07-23 & 2022-07-28 \\
SMB & 2022-07-26 & 2022-07-28 \\
\end{longtable}

As mentioned earlier, the data set contains a total of 484 columns, out
of which 426 are hydro acoustic signals or frequencies. The frequencies
go from F45 to F260, and they increase by 0.5. The other 58 columns
contain information such as species, unique identifier for the fish and
physical characteristics like sex, air bladder weight, tissue sample,
length, among others. It is possible that these measures where taken via
more aggressive methods; our analysis focuses on just the hydro acoustic
signals, as a form to explore non-lethal methods to deal with species
classification and problems adjacent to it.

For the purpose of this paper, we do not focus on the temporal data
aspect of this problem, though it is possible to start training models
based on this particular attribute of the data set. We leave this aspect
to be explored in future research.

\hypertarget{missing-values}{%
\subsection{Missing Values}\label{missing-values}}

\hypertarget{tbl-missing-values}{}
\begin{longtable}[]{@{}lrrr@{}}
\caption{\label{tbl-missing-values}Missing frequencies in the
dataset}\tabularnewline
\toprule\noalign{}
Species & F90-F170 & F173-F260 & Rows with any missing \\
\midrule\noalign{}
\endfirsthead
\toprule\noalign{}
Species & F90-F170 & F173-F260 & Rows with any missing \\
\midrule\noalign{}
\endhead
\bottomrule\noalign{}
\endlastfoot
LT & 2,964 & 183 & 3,147 \\
LW & 789 & 0 & 789 \\
SMB & 306 & 0 & 306 \\
Total & 4,059 & 183 & 4,242 \\
\end{longtable}

As observed in Table~\ref{tbl-missing-values}, our data set has 4,242
rows with missing values, most of which come from the Lake Trout species
which is also the species for which the most data was collected. It is
possible that some of these values are missing at random; for example,
the part of the machine that gathered those frequencies was not working
that day. While some imputation methods were considered to replace
missing values, given the small number of fish in the sample, the
consequences of imputing missing values are unknown, and specially when
focusing on dimensionality reduction we want to avoid introducing more
noise than there already is to perceive signals clearly. Consequently,
we drop any rows that have missing values in the frequencies (4,242 in
total).

After removing rows with missing frequencies, the resulting data set
contains 10,333 rows. Table~\ref{tbl-data-overview-clean} shows the
resulting data set. The number of LT fish is reduced by half. Based on
these results, the initial dimensionality reduction will be performed in
a binary classification setting between Lake Trout (LT) and Lake
Whitefish (LW), given that they are the classes with the most pings
(4,101 and 3,496), compared to Smallmouth Bass (2,736).

\hypertarget{tbl-data-overview-clean}{}
\begin{longtable}[]{@{}
  >{\raggedright\arraybackslash}p{(\columnwidth - 10\tabcolsep) * \real{0.1067}}
  >{\centering\arraybackslash}p{(\columnwidth - 10\tabcolsep) * \real{0.2133}}
  >{\centering\arraybackslash}p{(\columnwidth - 10\tabcolsep) * \real{0.1733}}
  >{\centering\arraybackslash}p{(\columnwidth - 10\tabcolsep) * \real{0.1867}}
  >{\centering\arraybackslash}p{(\columnwidth - 10\tabcolsep) * \real{0.1600}}
  >{\centering\arraybackslash}p{(\columnwidth - 10\tabcolsep) * \real{0.1600}}@{}}
\caption{\label{tbl-data-overview-clean}Overview of the dataset with no
missing values}\tabularnewline
\toprule\noalign{}
\begin{minipage}[b]{\linewidth}\raggedright
Species
\end{minipage} & \begin{minipage}[b]{\linewidth}\centering
Number of fish
\end{minipage} & \begin{minipage}[b]{\linewidth}\centering
Total pings
\end{minipage} & \begin{minipage}[b]{\linewidth}\centering
Median pings
\end{minipage} & \begin{minipage}[b]{\linewidth}\centering
Min. pings
\end{minipage} & \begin{minipage}[b]{\linewidth}\centering
Max. pings
\end{minipage} \\
\midrule\noalign{}
\endfirsthead
\toprule\noalign{}
\begin{minipage}[b]{\linewidth}\raggedright
Species
\end{minipage} & \begin{minipage}[b]{\linewidth}\centering
Number of fish
\end{minipage} & \begin{minipage}[b]{\linewidth}\centering
Total pings
\end{minipage} & \begin{minipage}[b]{\linewidth}\centering
Median pings
\end{minipage} & \begin{minipage}[b]{\linewidth}\centering
Min. pings
\end{minipage} & \begin{minipage}[b]{\linewidth}\centering
Max. pings
\end{minipage} \\
\midrule\noalign{}
\endhead
\bottomrule\noalign{}
\endlastfoot
LT & 12 & 4,101 & 189.5 & 16 & 1,379 \\
LW & 8 & 3,496 & 312.0 & 96 & 1,619 \\
SMB & 10 & 2,736 & 210.0 & 25 & 1,071 \\
\end{longtable}

\hypertarget{methods}{%
\section{Methods}\label{methods}}

In this section we compare several techniques for variable selection and
dimensionality reduction. We begin with regularization methods, which
fit a model containing all \(p\) predictors while constraining the
coefficient estimates, shrinking them towards zero. We use the standard
Lasso and a recent variant, the Whitening Lasso (WLasso), designed for
high-dimensional data with highly correlated predictors. We then
consider PCA, random forests and boosting.

\hypertarget{lasso-regularization}{%
\subsection{Lasso Regularization}\label{lasso-regularization}}

Lasso is a shrinkage method that adds an \(\ell_1\) penalty to the
model's loss function. When the tuning parameter \(\lambda\) is
sufficiently large, some coefficients are set exactly to zero, so Lasso
returns a sparse model that performs variable selection and is easier to
interpret. Because our response is binary (Lake trout vs.~Lake
whitefish), we use Lasso with logistic regression. Writing \(y_i = 1\)
for lake trout, \(x_i\) for the vector of standardized frequencies of
ping \(i\), and \(w_i\) for its weight, the coefficients are estimated
by

\[
\hat{\beta} = \arg\min_{\beta_0,\beta}\; -\frac{1}{n}\sum_{i=1}^{n} w_i\Big[y_i(\beta_0 + x_i^\top\beta) - \log\big(1 + e^{\beta_0 + x_i^\top\beta}\big)\Big] + \lambda \sum_{j=1}^{p}|\beta_j| .
\]

Larger values of \(\lambda\) shrink more coefficients to zero, while
\(\lambda = 0\) recovers ordinary logistic regression. The intercept is
not penalized.

\textbf{Preprocessing.} Each frequency was standardized to mean zero and
unit variance, as common practice when using Lasso regularization. The
means and standard deviations were computed on the training data only
and then applied to the test data, so that no information from the test
set would enter the model fitting. Additionally, classes differ in the
number of fish, and also each fish differs in its number of pings. To
prevent fish with many pings (and the larger class) from dominating the
fit, each ping received a weight inversely proportional to the number of
pings of its fish and to the number of fish in its class, rescaled to
average one. This scheme is know as weighted logistic regression, with
it each fish contributes equally within its class, and the two classes
contribute equally to each other. This helps prevent the model from
simply guessing the majority class to achieve a higher accuracy, instead
we penalize mistakes in classifying the lower represented class more to
account for the fact that the class is underrepresented.

\textbf{Data splitting and evaluation.} Pings from the same fish are not
independent, so all splitting was done by fish rather than by ping.
Splitting by ping would place pings of the same individual in both the
training and test data, which constitutes data leakage and overestimates
performance. After removing pings with missing values in the F90--F260
bands, 20 fish remained (12 lake trout and 8 lake whitefish), out of the
original 32. The number of pings per fish varies widely (from 16 to
1,619), so the number of pings in any subset cannot be controlled
directly. For this reason we report fish-level and balanced accuracy in
addition to ping-level accuracy.

First, 4 fish (2 lake trout and 2 lake whitefish) were randomly set
aside as a test set (seed = 16). They were not used for fitting or
tuning, and they are evaluated once, at the end. The same test fish are
used for every method in this study (Lasso, WLasso, random forest,
boosting), which allows a direct comparison of their performance on the
same unseen individuals.

The remaining 16 fish (10 lake trout and 6 lake whitefish) form the
development set. They were randomly divided into \(K=4\) folds of 4
fish, with species spread as evenly as possible so that each fold
contains both species. The penalty \(\lambda\) was chosen by \(K\)-fold
cross-validation: for each fold, the model was fitted on the other three
folds along a grid of \(\lambda\) values and the held-out fold was
scored, and the scores were averaged across folds. Folds were formed
from whole fish, so no individual appears in both the fitting and the
scoring data. We used the binomial deviance as the cross-validation
criterion, because, unlike accuracy, it evaluates the predicted
probabilities to gauge model confidence as opposed to accuracy. We
considered both the value of \(\lambda\) with the lowest cross-validated
error (\(\lambda_{\min}\)) and the largest value within one standard
error of the minimum (\(\lambda_{1\mathrm{SE}}\)). The final model was
then refitted on all 16 development fish with the chosen \(\lambda\) and
evaluated once on the 4 test fish.

\emph{The fish-level classification was obtained by averaging the
predicted probabilities of lake trout over a fish's pings, with a
threshold of 0.5. This is mainly because the fish is the real unit: for
every ping, the model outputs a probability that it came from a lake
trout. A single fish has hundreds or thousands of pings, so we get
hundreds or thousands of probabilities for that fish. To get one answer
per fish, we average them, and if the average is above 0.5 we call the
fish a lake trout, otherwise a whitefish.}

\textbf{Rationale.} A single train/validation/test split would leave
very few fish for each purpose, and cross-validation on the development
fish makes use of all 16 for both fitting and tuning. We chose \(K=4\)
so that every fold contains several fish of both species, while each
model is still fitted on 12 of the 16 fish. A fixed test set, identical
for all methods, was kept so that the final comparison between methods
is made on data that played no role in tuning or model selection.

\textbf{Limitations.} With 20 fish in total, all performance estimates
are uncertain. The test set contains only 4 fish, so fish-level accuracy
can only take the values 0, 25, 50, 75 or 100\%, and the number of pings
differs greatly between its two species. Test results should therefore
be interpreted together with the cross-validated results, not on their
own.

\hypertarget{results}{%
\section{Results}\label{results}}

\hypertarget{discussion}{%
\section{Discussion}\label{discussion}}

\hypertarget{sec-appendix-a}{%
\section{Appendix}\label{sec-appendix-a}}

\setcounter{figure}{0}
\renewcommand{\thefigure}{A\arabic{figure}}

\setcounter{table}{0}
\renewcommand{\thetable}{A\arabic{table}}



\end{document}
